% -----------------------------------------------------------
\section{Justification, Justify}
% -----------------------------------------------------------
	\label{JUSTIFICATION}
	
% % ----------------------
% \subsection{See also}
% % ----------------------

% ----------------------
\subsection{1828 American Dictionary of the English Language} % *1828
% ----------------------
	\label{JUSTIFICATION-1828}
	
	% -----
	\subsubsection{Justification}
	% -----

		\begin{compactitem}
			\item The act of justifying; a showing to be just or conformable to law, rectitude or propriety; vindication; defense. The court listened to the evidence and arguments in justification of the prisoner's conduct. Our disobedience to God's commands admits no justification
			\item Absolution.
			\item In law, the showing of a sufficient reason in court why a defendant did what he is called to answer. Pleas in justification must set forth some special matter.
			\item In theology, remission of sin and absolution from guilt and punishment; or an act of free grace by which God pardons the sinner and accepts him as righteous, on account of the atonement of Christ.
		\end{compactitem}
		
	% -----
	\subsubsection{Justify}
	% -----
	
		\begin{compactitem}
			\item To prove or show to be just, or conformable to law, right, justice, propriety or duty; to defend or maintain; to vindicate as right. We cannot justify disobedience or ingratitude to our Maker. We cannot justify insult or incivility to our fellow men. Intemperance, lewdness, profaneness and dueling are in no case to be justified.
			\item In theology, to pardon and clear from guilt; to absolve or acquit from guilt and merited punishment, and to accept as righteous on account of the merits of the Savior, or by the application of Christ's atonement to the offender.
			\item To cause another to appear comparatively righteous, or less guilty than one's self.
			\item To judge rightly of.
			\item To accept as just and treat with favor.
		\end{compactitem}

% % ----------------------
% \subsection{Oxford English Dictionary} % *OED
% % ----------------------
	% \label{JUSTIFICATION-OED}

	% \begin{compactitem}
		% \item[]
	% \end{compactitem}
	
% % ----------------------
% \subsection{Restoration Lexicon} % *RL *RESTORATION LEXICON
% % ----------------------
	% \label{JUSTIFICATION-OED}

	% \begin{compactitem}
		% \item[]
	% \end{compactitem}

% ----------------------
\subsection{Scriptural Usage} % *SCRIPTURES
% ----------------------
	\label{JUSTIFICATION-SCRIPTURES}

	% % -----
	% \subsubsection{Old Testament} % *OT
	% % -----
		% \label{JUSTIFICATION-OT}
	
		% % --
		% \paragraph{KJV}
		% % --
			% \label{JUSTIFICATION-OT-KJV}
	
			% \begin{tabular}{ll}
				% Total occurrences in the OT & 
			% \end{tabular}
			
			% \begin{compactdesc}
				% \item[]
			% \end{compactdesc}
			
		% % --
		% \paragraph{Hebrew Bible}
		% % --
			% \label{JUSTIFICATION-HEBREW}
		
			% %
			% \subparagraph{\texorpdfstring{\he{sample word}}{}} % paste actual hebrew text here
			% %
				% \label{PAGA} % whatever is in step bible without periods
			
				% \begin{compactdesc}
					% \item[HALOT , PDF ] \textbf{} % Use the 1994 PDF
						% \begin{compactitem}
							% \item[1]
						% \end{compactitem}
					% \item[BDB , PDF ] \textbf{}
						% \begin{compactitem}
							% \item[1]
						% \end{compactitem}
					% \item[DCH PDF ] \textbf{} % don't include the actual page number, they repeat from volume to volume
						% \begin{compactitem}
							% \item[1]
						% \end{compactitem}
				% \end{compactdesc}
				
				% \begin{compactdesc}
					% \item[Some OT uses] \textbf{}
						% \begin{compactdesc}
							% \item[]
						% \end{compactdesc}
				% \end{compactdesc}
			
		% % --
		% \paragraph{Septuagint}
		% % --
			% \label{JUSTIFICATION-LXX}
		
			% %
			% \subparagraph{\texorpdfstring{\gr{sample word}}{}}
			% %
		
	% -----
	\subsubsection{New Testament} % *NT
	% -----
		\label{JUSTIFICATION-NT}
	
		% --
		\paragraph{KJV}
		% --
			\label{JUSTIFICATION-NT-KJV}		
			
			%
			\subparagraph{Justification}
			%
	
				\begin{tabular}{ll}
					Total occurrences in the NT & 3
				\end{tabular}
				
				\begin{compactdesc}
					\item[{\hyperref[ROMANS-3]{Romans 3}}] \phantom{.}
						\begin{compactitem}
							\item Various verses in this chapter.
						\end{compactitem}
					\item[{\hyperref[ROMANS-4-25]{Romans 4:25}}]
					\item[{\hyperref[ROMANS-5-16]{Romans 5:16}}]
					\item[{\hyperref[ROMANS-5-18]{Romans 5:18}}]
				\end{compactdesc}
				
			%
			\subparagraph{Justify}
			%
	
				\begin{tabular}{ll}
					Total occurrences in the NT & 37
				\end{tabular}
				
				\begin{compactdesc}
					\item[{\hyperref[LUKE-16-15]{Luke 16:15}}]
					\item[{\hyperref[GALATIANS-2-16]{Galatians 2:16--21}}]
				\end{compactdesc}
			
		% --
		\paragraph{Greek New Testament} % *GREEK
		% --
			\label{JUSTIFICATION-GREEK}
			
			%
			\subparagraph{\texorpdfstring{\gr{δικαιόω}}{}}
			%
				\label{DIKAIOO}
			
				\begin{compactdesc}
					\item[BDAG 220b] \textbf{}
						\begin{compactitem}
							\item[1] \textbf{to take up a legal cause, \textit{show justice, do justice, take up a cause}}
							\item[2] \textbf{to render a favorable verdict, \textit{vindicate}}
							\item[2.a] as activity of humans \textit{justify, vindicate, treat as just}
							\item[2.b] of experience or activity of transcendent figures, especially in relation to humans
							\item[2.b.$\alpha$] of wisdom
							\item[2.b.$\beta$] of God \textit{be found in the right, be free of charges}...Paul, who has influenced later writers...uses the word almost exclusively of God's judgment. As affirmative verdict, \hyperref[ROMANS-2-13]{Romans 2:13}. Especially of persons \gr{δικαιοῦσθαι} \textit{be acquitted, be pronounced as treated as righteous} and thereby become \gr{δίκαιος}, receive the divine gift of \gr{δικαιοσύνη} through faith in Christ Jesus and apart from \gr{νόμος} as a basis for evaluation...since Paul views God's justifying action in close connection with the power of Christ's resurrection, there is sometimes no clear distinction between the justifying action of acquittal and the gift of new life through the Holy Spirit as God's activity in promoting uprightness in believers. Passages of this nature include \hyperref[ROMANS-3-26]{Romans 3:26, 30}, \hyperref[ROMANS-4-5]{4:5} (on \gr{δικαιοῦν τὸν ἀσεβῆ} compare the warning against accepting \gr{δῶρα} to arrange acquittal [\hyperref[EXODUS-23-7]{Exodus 23:7} and \hyperref[ISAIAH-5-23]{Isaiah 5:23}]; \gr{δικαιούμενοι δωρεάν} [\hyperref[ROMANS-3-24]{Romans 3:24} is therefore all the more pointed); \hyperref[ROMANS-8-30]{8:30, 33} (\hyperref[ISAIAH-50-8]{Isaiah 50:8}); \hyperref[GALATIANS-3-8]{Galatians 3:8}.
							\item[3] \textbf{to cause someone to be released from personal or institutional claims that are no longer to be considered pertinent or valid, \textit{make free/pure}}... in our literature passive \gr{δικαιοῦμαι} \textit{be set free, made pure} \gr{ἀπό} from...\gr{ἀπὸ πάντων ὧν οὐκ ἠδυνήθητε ἐν νόμω Μωϋσέως δικαιωθῆναι} \textit{from everything from which you could not be freed by the law of Moses} (\hyperref[ACTS-13-38]{Acts 13:38}); compare \hyperref[ACTS-13-39]{verse 39} \gr{ὁ ἀποθανὼν δεδικαίωται ἀπὸ τ. ἁμαρτίας} \textit{the one who died is freed from sin} (\hyperref[ROMANS-6-7]{Romans 6:7})...In the context of \hyperref[1CORINTHIANS-6-11]{1 Corinthians 6:11} \gr{ἐδικαιώθητε} means \textit{you have become pure}.---In the language of the mystery religions...\gr{δικαιοῦσθαι} refers to a radical inner change which the initiate experiences...and approaches the sense `become deified'. Some are inclined to find in \hyperref[1TIMOTHY-3-16]{1 Timothy 3:16} a similar use; but see under \#4.
							\item[4] \textbf{to demonstrate to be morally right, \textit{prove to be right}}, passive of God \textit{is proved to be right} (\hyperref[ROMANS-3-4]{Romans 3:4}).
						\end{compactitem}
					% \item[CGL , PDF ] \textbf{}
						% \begin{compactitem}
							% \item 
						% \end{compactitem}
					% \item[LSJ , PDF ] \textbf{}
						% \begin{compactitem}
							% \item 
						% \end{compactitem}
				\end{compactdesc}
				
				% \begin{tabular}{ll}
					% Total occurrences in the NT & 
				% \end{tabular}
				
				\begin{compactdesc}
					\item[Some NT uses] \textbf{}
						\begin{compactdesc}
							\item[{\hyperref[LUKE-16-15]{Luke 16:15}}]
						\end{compactdesc}
				\end{compactdesc}
				
			% %
			% \subparagraph{TDNT: \texorpdfstring{\gr{sample word}}{}  (3:536--556)}
			% %
				% \label{KATALLAGE-TDNT}
		
			%
			\subparagraph{\texorpdfstring{\gr{δικαίωσις}}{}}
			%
				\label{DIKAIOSIS}
			
				\begin{compactdesc}
					\item[BDAG 221b] \textbf{}
						\begin{compactitem}
							\item \textbf{\textit{justification, vindication, acquittal}}...as a process as well as its result
								\begin{compactitem}
									\item This entry glosses \hyperref[ROMANS-5-18]{Romans 5:18} as ``\textit{acquittal that brings life}.''
								\end{compactitem}
						\end{compactitem}
					% \item[CGL , PDF ] \textbf{}
						% \begin{compactitem}
							% \item 
						% \end{compactitem}
					% \item[LSJ , PDF ] \textbf{}
						% \begin{compactitem}
							% \item 
						% \end{compactitem}
				\end{compactdesc}
				
				% \begin{tabular}{ll}
					% Total occurrences in the NT & 
				% \end{tabular}
				
				% \begin{compactdesc}
					% \item[Some NT uses] \textbf{}
						% \begin{compactdesc}
							% \item[]
						% \end{compactdesc}
				% \end{compactdesc}
				
			% %
			% \subparagraph{TDNT: \texorpdfstring{\gr{sample word}}{}  (3:536--556)}
			% %
				% \label{KATALLAGE-TDNT}
		
	% % -----
	% \subsubsection{Book of Mormon} % *BOM
	% % -----
		% \label{JUSTIFICATION-BOM}
	
		% \begin{tabular}{ll}
			% Total occurrences in the BOM & 
		% \end{tabular}
		
		% \begin{compactdesc}
			% \item[]
		% \end{compactdesc}
		
	% -----
	\subsubsection{Doctrine and Covenants} % *DC
	% -----
		\label{JUSTIFICATION-DC}
	
		% \begin{tabular}{ll}
			% Total occurrences in the DC & 
		% \end{tabular}
		
		\begin{compactdesc}
			\item[{\hyperref[DC98]{DC 98}}] The words ``justification'' and ``justify'' appear a number of times in this section; they seem to have a more strongly \textit{legal} sense here, as opposed to a theological sense with a legal flavor, than in other places.
		\end{compactdesc}
		
	% % -----
	% \subsubsection{Pearl of Great Price} % *PGP
	% % -----
		% \label{JUSTIFICATION-PGP}
	
		% \begin{tabular}{ll}
			% Total occurrences in the PGP & 
		% \end{tabular}
		
		% \begin{compactdesc}
			% \item[]
		% \end{compactdesc}
		
% % ----------------------
% \subsection{Key Phrases} % *KEYPHRASES *PHRASES
% % ----------------------
	% \label{JUSTIFICATION-PHRASES}

	% \begin{compactdesc}
		% \item[] \textbf{\#times} | list
			% % \begin{compactdesc}
				% % \item[Bible anchor verses] \phantom{.}
					% % \begin{compactdesc}
						% % \item[Romans 5:5] \gr{}
					% % \end{compactdesc}
				% % \item[Commentary] ---
			% % \end{compactdesc}
	% \end{compactdesc}
	
% % ----------------------
% \subsection{Bible Dictionaries} % *DICTIONARIES
% % ----------------------
	% \label{JUSTIFICATION-BD}

% % ----------------------
% \subsection{Restoration Usage} % *RESTORATION
% % ----------------------
	% \label{JUSTIFICATION-RESTORATION}

	% % -----
	% \subsubsection{Joseph Smith} % *JS
	% % -----
		% \label{JUSTIFICATION-JS}
	
		% \begin{compactdesc}
			% \item[{\hyperref[KEY]{Date/Source}}]
		% \end{compactdesc}
		
	% % -----
	% \subsubsection{Temple Ordinances}
	% % -----
		% \label{JUSTIFICATION-TEMPLE}
	
		% \begin{compactdesc}
			% \item[{\hyperref[KEY]{Date/Source}}]
		% \end{compactdesc}
	
	% % -----
	% \subsubsection{Brigham Young} % *BY
	% % -----
		% \label{JUSTIFICATION-BY}
	
		% \begin{compactdesc}
			% \item[{\hyperref[KEY]{Date/Source}}]
		% \end{compactdesc}
	
% % ----------------------
% \subsection{Commentary and Studies} % *COMMENTARY *STUDIES
% % ----------------------
	% \label{JUSTIFICATION-COMMENTARY}

	% % -----
	% \subsubsection{Key Points}
	% % -----
		% \label{JUSTIFICATION-KEYS}
	
		% \begin{compactitem}
			% \item
		% \end{compactitem}
		
	% % -----
	% \subsubsection{Sample Study}
	% % -----
		% \label{JUSTIFICATION-study}